%% Section 2: Related Work
\section{Related Work}
\label{sec:related}
The pursuit of epistemic fortitude is situated at the intersection of several active research areas in large language model (LLM) development: AI alignment, agent-level self-correction, and multi-agent reasoning. This section surveys the foundational and state-of-the-art literature in these domains to establish the context for our proposed architecture.


\subsection{The Sycophancy Problem in Aligned Language Models}
\label{subsec:related-sycophancy}

The dominant paradigm for aligning LLMs with human values is Reinforcement Learning from Human Feedback (RLHF), a technique detailed by \citet{ouyang2022training} that trains models to produce outputs that human labelers prefer. While highly effective at improving helpfulness and safety, this process has a critical side effect: it inadvertently incentivizes \textit{sycophancy}, the tendency for a model to align with a user's stated beliefs at the expense of factual accuracy.

\citet{sharma2023towards} provided a foundational analysis of this phenomenon, showing that human preference judgments systematically favor agreeable or flattering responses. This creates a perverse incentive where models learn that agreeableness is a reliable proxy for high reward, leading them to prioritize user approval over factual integrity. A comprehensive survey by \citet{cheng2025sycophancy} identifies several manifestations, including \textit{opinion sycophancy} (mirroring user beliefs), \textit{mimicry of user errors}, and \textit{feedback sycophancy}, where a model wrongly admits a mistake simply because a user challenges it.

Recent interpretability research by \citet{wynn2025disentangling} suggests that sycophantic agreement and genuine, truth-based agreement correspond to distinct and independently controllable directions in a model's activation space, implying that sycophancy is a targetable mechanism rather than an unavoidable trade-off. \citet{zhang2025pressure} propose \textit{Pressure-Tune}, a method that fine-tunes models on synthetic adversarial dialogues with chain-of-thought rationales that explicitly reject user misinformation. \citet{zhao2025kbqg} similarly target RLHF's sycophancy and confirmation biases through a bias-corrected reinforcement learning framework (BC-RLHF), replacing raw human preference signals with a more carefully calibrated reward design for knowledge base question generation. However, training-based approaches have shown limited success in fully eliminating sycophancy, as the fundamental tension between user satisfaction and factual accuracy persists within monolithic architectures. This limitation motivates architectural solutions that externalize epistemic oversight.

The real-world consequences of sycophancy extend beyond theoretical concern. \citet{bo2026invisible} demonstrate that sycophantic LLMs actively mislead novice users during problem-solving tasks, producing measurable harm to learning outcomes---an ``invisible saboteur'' effect where users trust incorrect AI feedback precisely because it agrees with their misconceptions. In software engineering, where LLMs are increasingly embedded in development workflows~\citep{belozerov2025llms}, such failure modes pose particular risks: a code review assistant that caves to a developer's incorrect hypothesis can propagate bugs into production systems.


\subsection{Evaluating Sycophancy: SycEval and Related Benchmarks}
\label{subsec:related-syceval}

Concurrent with mitigation efforts, researchers have developed evaluation frameworks to systematically measure sycophantic behavior. \citet{fanous2025syceval} introduce \textit{SycEval}, a benchmark for evaluating LLM sycophancy across ChatGPT-4o, Claude-Sonnet, and Gemini-1.5-Pro. SycEval provides standardized test scenarios and scoring criteria that quantify how readily models abandon correct positions under user pressure.

Our work is complementary to but distinct from SycEval in two respects. First, SycEval focuses on \textit{evaluating} sycophancy as a model property, whereas we focus on \textit{intervening} to prevent it through architectural means. Second, SycEval measures sycophancy within a single model, while our framework introduces a multi-agent architecture that modifies the model's behavior at inference time without retraining. The evaluation methodology we develop (Section~\ref{sec:experiments}) shares SycEval's emphasis on systematic scoring but extends it to measure the effectiveness of architectural interventions across multiple model families.


\subsection{Agent-Level Self-Correction: Capabilities and Limitations}
\label{subsec:related-self-correction}

A significant body of research has explored methods to enable a single LLM to improve its own outputs. Frameworks like \textit{Self-Refine} \citep{madaan2023self} allow a model to iteratively generate feedback on its own draft and refine it, improving output quality without additional training. Similarly, the \textit{Self-Taught Reasoner} (STaR) framework \citep{zelikman2022star} enables a model to bootstrap its reasoning capabilities by fine-tuning on reasoning chains that lead to correct answers.

However, the efficacy of intrinsic self-correction is fundamentally limited. A critical line of inquiry has revealed a ``dark side'' to this capability: without an external ground-truth oracle, self-correction prompts can trigger \textbf{answer wavering}, causing a model to change a correct initial answer to an incorrect one \citep{zhang2025understanding}. \citet{kamoi2024can} provide a comprehensive analysis showing that self-correction often fails when models lack access to external verification.

The problem is compounded when handling external user feedback. \citet{jiang2025feedback} identified \textit{Feedback Friction}, a phenomenon where LLMs consistently resist incorporating even high-quality, correct external feedback due to high internal confidence in their initial answer. Conversely, \citet{liu2025user} showed that real-world user feedback is often a ``noisy'' and ambiguous learning signal, making it difficult for an agent to distinguish between a valid correction and a misleading statement.

This triad of limitations---unreliable intrinsic correction, resistance to valid external correction, and vulnerability to noisy feedback---establishes the need for external supervisory mechanisms that can evaluate the substance of disagreements rather than relying on the model's own judgment.


\subsection{Multi-Agent Systems for Enhanced Reasoning}
\label{subsec:related-multiagent}

To overcome the limitations of a single agent, researchers have turned to multi-agent systems. The \textit{Multi-Agent Debate} (MAD) framework, introduced by \citet{du2023improving}, showed that having multiple LLM instances propose, critique, and refine solutions can significantly improve performance on complex reasoning tasks and reduce hallucinations. This paradigm has been extended with persistent memory (MeMAD) \citep{ling2025memad} and frameworks that enforce cognitive diversity by assigning different reasoning strategies to each agent (DMAD) \citep{anonymous2025dmad}. More recent work demonstrates multi-agent specialization in knowledge-intensive domains: \citet{meier2025structured} orchestrate multiple smaller specialized models through structured workflows---including a Delphi consensus mechanism for iterative refinement---for causal discovery, while \citet{xiong2025delphiagent} propose DelphiAgent, a fact verification framework where agents with distinct roles make independent factuality judgments and converge through iterative consensus rounds.

Despite these successes, the assumption that more debate is always better is flawed. \citet{wynn2025talk} provided a crucial empirical demonstration that multi-agent debate can be actively harmful. The primary failure mode is a form of \textit{multi-agent sycophancy}, where agents frequently abandon their own correct reasoning to agree with a peer's flawed argument, prioritizing consensus over correctness. This dynamic reveals that if all agents share the same underlying biases instilled by RLHF, debate can become an echo chamber that reinforces error.

These findings suggest that peer-based debate alone may be insufficient, and that hierarchical architectures with differentiated agent roles and explicit supervisory authority offer a promising alternative direction---one that we explore in this work.


\subsection{Constitutional AI as a Principled Supervisory Framework}
\label{subsec:related-constitutional}

The design of our Arbiter Agent draws directly from the \textit{Constitutional AI} (CAI) framework developed by \citet{bai2022constitutional}. CAI guides model behavior through a ``constitution''---a set of explicit, human-written principles---rather than relying on direct human preference labels. The training process involves a supervised phase where the model critiques and revises its own outputs according to the constitution, followed by reinforcement learning using AI-generated feedback (RLAIF) based on those same principles.

The primary advantage of CAI is that it replaces the sycophancy-inducing signal of raw human preference with a scalable, transparent, and principle-driven AI feedback signal \citep{anthropic2023constitutional}. However, when implemented within a single model, Constitutional AI faces several limitations in conversational contexts. Helpfulness and accuracy often conflict during user contradictions, and CAI provides no principled way to prioritize when these objectives are encoded in a single reward signal. All principles are collapsed into a single optimization target during training, making fine-grained trade-offs difficult. Furthermore, constitutional principles are fixed at training time and cannot adapt to context-specific epistemic requirements at runtime.

These limitations suggest that constitutional principles may be more effective when specialized across multiple agents rather than collapsed into a single model.


\subsection{Positioning This Work}
\label{subsec:related-positioning}

Our work synthesizes insights from these research threads. From sycophancy research, we adopt the problem formulation and extend it with an intervention-focused approach, complementing evaluation frameworks like SycEval~\citep{fanous2025syceval} with an architectural solution. From self-correction literature, we learn the limitations of single-agent epistemic oversight and the need for external anchoring mechanisms. From multi-agent systems, we borrow the architectural pattern but modify it with hierarchy to prevent false consensus---distinguishing our Primary-Arbiter design from consensus-based approaches such as DelphiAgent~\citep{xiong2025delphiagent}. From Constitutional AI, we derive the principle-based instruction design for the Arbiter, extending CAI's application to conversational epistemic integrity.

The novel contribution is the integration: a multi-agent architecture where epistemic principles guide a specialized fact-checking agent, activated through hybrid routing when conversational context demands factual adjudication. Critically, we validate this architecture across four model families and provide ablation evidence that the architectural separation itself---not merely the prompt content---drives a substantial portion of the improvement.
